\section{线性表示、Superposition 与 Sparse Autoencoder}

\subsection{Features 与 Circuits：理解的基本单元}

Inception V1 中许多神经元有清晰的可解释含义：曲线、汽车、车轮、狗耳朵。Chris Olah 展示了一个完整的汽车检测器电路：连接到窗户检测器（上方）、车轮检测器（下方）、车身检测器（中间）——这\textbf{就是}检测汽车的算法，直接从权重中读出来。

但问题是：不是所有神经元都可解释。\textbf{多义神经元}（polysemantic neurons）响应多个不相关的事物。这带来指数级的复杂性：

\begin{warningbox}{多义性导致指数爆炸}
如果两个多义神经元各响应 3 个概念，它们之间的权重就有 $3 \times 3 = 9$ 种可能的交互需要考虑。更深层的问题是：高维空间具有指数级的体积——如果不能将其分解为独立部分，理解的复杂度将无法控制。

\textbf{单义特征}（monosemantic features）——具有单一清晰含义的特征——允许独立推理，从而避免指数爆炸。这就是为什么追求单义性如此重要。
\end{warningbox}

\subsection{线性表示假说}

\begin{importantbox}{方向有意义}
激活空间中的\textbf{方向}具有语义含义。更多的激活 = 更高的检测置信度。

Word2Vec 的经典例子：King - Man + Woman = Queen；Sushi - Japan + Italy = Pizza。能够加减向量并得到有意义的结果，反映了线性结构。

\textit{``这实际上就是正在发生的根本性的事情——方向有意义。''} 到目前为止在自然神经网络中观察到的一切都与线性表示假说一致。
\end{importantbox}

Chris 也提到一些关于非线性表示的新工作（多维特征/流形），但他对线性假说的态度很务实——引用了一个精彩的类比：

\begin{knowledgebox}{热质理论的教训}
即使是错误的理论（如热质理论），如果认真对待并推到极致，也能产生实际成果——\textbf{燃烧机就是由相信热质理论的人开发的}。

\textit{``认真对待假说并将其推到极限是有价值的。''} 线性表示假说即使最终被证明不完全正确，在其有效范围内推进也能产生重大发现。
\end{knowledgebox}

\subsection{Superposition 假说}

500 维的词嵌入不可能只对应 500 个概念——模型需要表示远多于维度数的概念。解决方案：利用高维空间的几何特性 + 概念的稀疏性（``日本''和``意大利''很少同时出现）。

\begin{importantbox}{``更大稀疏网络的影子''}
\textit{``神经网络可能是更大、更稀疏的神经网络的影子，我们看到的是这些投影。''}

想象一个``楼上模型''——巨大但稀疏，每个神经元都可解释。实际的神经网络是这个楼上模型的压缩投影。学习 = 构建楼上模型的高效压缩。梯度下降可能在秘密搜索极度稀疏模型的空间，然后将它们折叠成密集矩阵。

概念数量：Johnson-Lindenstrauss 引理表明，可嵌入的近似正交方向数是维度数的\textbf{指数级}。稀疏性和相关结构进一步增加了这个数字。
\end{importantbox}

\subsection{Sparse Autoencoder 突破}

如果 Superposition 假说正确，Dictionary Learning / Sparse Autoencoder 就是自然的解决方案。

\textbf{2023 年 10 月} ``Towards Monosemanticity''：在单层模型上使用 Sparse Autoencoder，\textit{``美丽的可解释特征就这样自然涌现出来''} 。发现了：阿拉伯语特征、希伯来语特征、Base64 特征、编程语言特征、特定上下文中的``the''（出现在数学上下文时预测``vector''和``matrix''）、Unicode 半字符交替模式。

训练两次模型，在两者中都能找到类似特征——Universality 再次被确认。\textit{``一个非常自然的技术就这样有效了……这实际上是一个非常好的状况。''}

\textbf{2024 年 5 月} 扩展到 Claude 3 Sonnet（生产模型）。需要\textbf{大量} GPU；Tom Henighan 的 Scaling Laws for Interpretability 帮助预测最优的 Sparse Autoencoder 大小和训练 token 数。

\begin{warningbox}{``令人毛骨悚然''的抽象多模态特征}
在 Claude 3 Sonnet 中发现了令人着迷但也令人不安的特征：

\begin{itemize}
    \item \textbf{安全漏洞特征}：同时响应``disable SSL''文本\textbf{和}点击 Chrome SSL 警告的图片
    \item \textbf{后门特征}：同时响应代码后门\textbf{和}隐藏摄像头设备的图片
    \item \textbf{欺骗/说谎特征}：强制激活后 Claude 开始说谎
    \item \textbf{寻求权力、发动政变、隐瞒信息}等特征
\end{itemize}

\textit{``这展示了这些概念有多抽象。''} 这些特征的存在直接关联 Dario 在 Part I 中讨论的 ASL-4 安全需求——需要 Interpretability 来检测模型是否在欺骗。
\end{warningbox}

\subsection{自动化可解释性的局限}

一个自然的想法是用 Claude 来标注 Sparse Autoencoder 发现的特征。Chris 承认自动化可解释性有其价值，但他\textit{``对此有些怀疑''}：

AI 给出的标签\textit{``在某种意义上是真的，但并没有真正抓住具体特征''}。这类似于数学家对计算机自动化证明的怀疑——你相信结论是正确的，但你没有\textbf{理解}。

更深层的安全顾虑：Ken Thompson 的经典论文``Reflections on Trusting Trust''——如果你用 AI 来验证 AI 的安全性，你能信任审计者吗？目前这不是大问题，但随着系统变得更强大，这个``谁来守卫守卫者''的问题将变得关键。

\subsection{未来方向：从微观到宏观}

\textbf{从特征到电路}：当前理解了表征（什么被激活），但还没理解计算过程（信息如何流动和变换）。干扰权重（Superposition 的人工产物）使电路分析更困难。

\textbf{``暗物质''问题}：当前的 Sparse Autoencoder 只能看到神经网络``物质''的一小部分。\textit{``就像早期天文学：随着我们建造更好的 Sparse Autoencoder……我们看到越来越多的星星。''}

\textbf{从微生物学到解剖学}：当前的 Mech Interp 是``神经网络的微生物学''——非常精细。但我们关心的问题是宏观的。需要攀升层次：分子生物学 $\to$ 细胞生物学 $\to$ 组织学 $\to$ 解剖学 $\to$ 动物学 $\to$ 生态学。或物理学的类比：个体粒子 $\to$ 统计物理 $\to$ 热力学。

\begin{knowledgebox}{神经网络有``器官''吗？}
\textit{``我希望存在比特征和电路大得多的东西。''} 是否存在类似心脏、大脑区域、呼吸系统的宏观结构？

不能直接跳到宏观——需要先理解微观结构，再研究连接模式。就像你不能跳过分子生物学直接研究解剖学。
\end{knowledgebox}

\subsection{人工 vs 生物神经网络}

Chris 认为神经科学家的工作``困难得多''。研究人工神经网络的优势清单令人印象深刻：

\begin{itemize}
    \item 记录\textbf{所有}神经元（不仅仅是能触及的）
    \item 输入任意数据
    \item 神经元在研究期间不变化
    \item 可以消融（删除）任何神经元
    \item 可以编辑任意连接权重
    \item 可以撤销所有更改
    \item 可以强制激活任何神经元到任意值
    \item 知道完整连接组，\textbf{含精确权重}（不仅是二值连接）
    \item 可以计算梯度
\end{itemize}

\textit{``我们拥有这么多超越神经科学家的优势，然后即使拥有所有这些优势，这仍然非常困难。如果对我们来说这么难，在神经科学的约束下似乎几乎不可能。''}

这也是为什么 Chris 积极从神经科学领域招人——\textit{``一个更容易的问题，但仍然非常困难''}。

\subsection{安全与美：Mech Interp 的双重召唤}

有人对神经网络``感到失望''：\textit{``只是简单规则的放大。''} Chris 的反驳极为优美：进化也``只是简单规则''——随机突变加自然选择。但它产生了生物学的全部壮丽。

\textit{``美在于简单性产生了复杂性。''} 神经网络内部创造了\textit{``人们通常不去看的巨大的复杂性和美。一个等待被发现的、令人难以置信的丰富结构。''}

最终的画面：电路向着损失函数的光芒生长——\textit{``这个我们培育出的有机体，而我们不知道我们培育出了什么。''} 这既是 Mech Interp 的科学动力，也是它的美学动力。

用 Alan Watts 的话收尾：\textit{``The only way to make sense out of change is to plunge into it, move with it, and join the dance.''}

\subsection{本章小结}

线性表示假说和 Superposition 假说构成了当前 Mech Interp 的理论基础。Sparse Autoencoder 从 Superposition 中提取可解释特征——从单层模型到 Claude 3 Sonnet 的成功扩展证明了方法的可行性。自动化可解释性有用但不能替代人类理解。未来的核心挑战是从``微生物学''走向``解剖学''——发现神经网络的宏观结构。Mech Interp 不仅服务于安全，也是对``我们到底创造了什么''这一深刻科学问题的探索。

